\section{The Learner}\label{sec:learner}

The learner is the Python package \code{gofer\_train}, built on
PyTorch~\citep{paszke2019pytorch}. It reads shards, trains one candidate per cycle, and
exports it. The v3 command-line tools keep their names and flags and now call into it.

\subsection{Data path}\label{sec:learner-data}

\paragraph{Loading.}
A source may be a legacy JSONL file, one shard, or a directory of shards. Directories are
ordered oldest to newest by modification time. Given a window of $W$ rows (from the
orchestrator, eq.~\eqref{eq:window}), the loader walks the sources from newest to oldest,
stops once $W$ rows are collected, and keeps exactly the newest $W$. Old shards are never
decompressed. Game identifiers are re-based per source, so games stay distinct across
shards.

\paragraph{Legacy JSONL.}
Old data is converted into shard conventions on load:
\begin{itemize}
  \item Ownership, stored with Black $=+1$, is multiplied by $-1$ on White-to-move rows.
  \item The field named \code{policy\_opp} is dropped, because in v3 it held the policy of
        the move \emph{before} the current position. The correct reply target
        \code{policy\_next}, added to JSONL in v4, is used when present.
  \item Game boundaries are recovered from resets of the move number, since v3 recorded
        no game id.
\end{itemize}
A test compares converted ownership with the raw file for every row.

\paragraph{Device residency.}
The window fits comfortably in accelerator memory. 50{,}000 $9\times9$ rows take
$50{,}000\times8\times81$ bytes $\approx$ 32\,MB of \code{uint8} planes, plus float
targets. So the whole window is copied to the training device once. Each step gathers a
batch by index on the device and converts planes to float there. This removes the v3
\code{DataLoader}, which built a tensor from a Python list for every row, every epoch, and
it removes host-to-device copies from the inner loop.

\paragraph{Game-level split.}
Positions from one game are strongly correlated.
\citet{silver2016mastering} sampled one position per game for exactly this reason when
training AlphaGo's value network. v3 split train and validation by row, so a validation
position usually had near-duplicates in training, and validation loss partly measured
memorisation. v4 assigns whole games to one side. A test checks that the two sets of game
ids are disjoint. One consequence: v4 validation losses are \emph{not comparable} with v3
numbers.

\paragraph{Recency sampling.}
With the default decay $\lambda=0$, batches are drawn uniformly from the window, as in
\katago{}. With $\lambda>0$, row $i$ of an $n$-row window (0 = oldest) is drawn with
probability proportional to $\exp(\lambda(i/n-1))$, so the newest data is sampled
$e^{\lambda}$ times as often as the oldest. \katago{} does not use this. It is offered as
a knob for small runs where the window is dominated by early, weaker games. The CPU preset
uses $\lambda=1.5$ and the GPU preset $\lambda=2$. Neither value has been tuned.

\paragraph{Symmetry augmentation.}
Each sample in a batch gets an independent random element of the dihedral group $D_4$.
The same element is applied to the spatial planes (all are point-local, including ko and
move history), to the $S^2$ board entries of both policy targets (the pass entry is left
alone), and to the ownership map. A test places one asymmetric stone together with its
policy mass and ownership, and checks two things: under all sampled symmetries the three
stay on the same point, and all eight images of the point are reached.

\subsection{Loss}\label{sec:learner-loss}\label{sec:learner-masking}

Let a batch $B$ have rows $i$ with full-search flag $f_i$, value target $z_i\in\{-1,0,1\}$,
policy target $\pi_i$, reply target $\pi^{\mathrm{opp}}_i$ (possibly empty), ownership
$o_i\in\{-1,0,1\}^{S^2}$ and final margin $s_i$ (possibly unknown). All targets are from
the side to move. Define the row sets
\begin{equation*}
  F = \{i : f_i = 1,\ \textstyle\sum\pi_i > 0\},\quad
  F^{\mathrm{opp}} = \{i : f_i = 1,\ \textstyle\sum\pi^{\mathrm{opp}}_i > 0\},\quad
  K = \{i : s_i \text{ known}\}.
\end{equation*}
With $\mathrm{CE}(p,q) = -\sum_m p(m)\log q(m)$ and $\mathrm{mean}_X$ the mean over a set
$X$ (zero if $X$ is empty), the loss is
\begin{align}
  L ={}& w_p\,\mathrm{mean}_{F}\,\mathrm{CE}(\pi_i,\hat\pi_i)
       + w_v\,\mathrm{mean}_{B}\,(z_i-\hat v_i)^2
       + w_o\,\mathrm{mean}_{B}\,\tfrac{1}{S^2}\lVert o_i-\hat o_i\rVert^2 \nonumber\\
      &+ w_{\mathrm{opp}}\,\mathrm{mean}_{F^{\mathrm{opp}}}\,\mathrm{CE}(\pi^{\mathrm{opp}}_i,\hat\pi^{\mathrm{opp}}_i)
       + w_s\,\mathrm{mean}_{K}\,H_1\!\left(\hat s_i - s_i/20\right)
       + w_d\,L_{\mathrm{distill}},
  \label{eq:loss}
\end{align}
where $H_1$ is the Huber loss~\citep{huber1964robust} with threshold 1, and
$\hat v,\hat o$ are $\tanh$ outputs. The default weights are $w_p=1$, $w_v=1.5$ (matching
\katago{}'s $c_g$), $w_o=0.15$, $w_{\mathrm{opp}}=0.15$ (matching \katago{}), $w_s=0.05$,
and $w_d=0$ unless a teacher is given. The auxiliary terms act as multi-task
regularisers~\citep{caruana1997multitask} for the shared trunk. Weight decay is decoupled from the
loss~\citep{loshchilov2019decoupled} and applied only to convolution and linear weights,
never to biases or batch-norm parameters~\citep{ioffe2015batch}.

\paragraph{How this differs from \katago{}.}
\begin{itemize}
  \item \textbf{Value.} Like AlphaGo Zero~\citep{silver2017mastering} and unlike \katago{}'s
        win/loss cross-entropy, the value is a squared error on a $\tanh$ output. The
        engine's ONNX contract is a single scalar in $[-1,1]$, and we kept it.
  \item \textbf{Ownership and score.} \katago{} uses cross-entropy for ownership and a
        pdf/cdf pair over a discretised score distribution. We use regression: squared
        error for ownership and Huber for score. This is simpler, but it loses the
        score-distribution information \katago{} exploits.
\end{itemize}
Both are simplifications, not claimed improvements.

\paragraph{Masking fast-search rows: a deliberate change with a known risk.}
\katago{} records only full-search turns. v4 records every turn and masks only the policy
term (sets $F$, $F^{\mathrm{opp}}$). The argument for this is that the outcome, ownership
and final score of a fast-search turn are exactly as valid as those of a full-search turn,
because the game produces them rather than the search. With $p=0.20$, each game then
yields about $1/p = 5\times$ more value, ownership and score rows at no extra self-play
cost.

The argument against it is the correlation that \citet{silver2016mastering} warned about.
Extra rows from the \emph{same} games add little independent information about $z$, and
they can overfit the value head to individual games. \katago{}'s choice may have been made
for exactly this reason.

We therefore treat masking as a hypothesis, not a result. It needs the ablation in
\cref{sec:conclusion}: identical runs with fast rows kept and discarded, compared in Elo
per unit of self-play compute. Game-level validation (\cref{sec:learner-data}) is what
makes the overfitting risk visible in the first place. A row-level split would hide it.

\subsection{Optimisation}\label{sec:learner-opt}

Each cycle warm-starts from the champion's weights (\code{--init-from}) with a fresh
optimiser. The optimiser options are:
\begin{itemize}
  \item \textbf{Optimiser.} SGD with Nesterov momentum 0.9~\citep{sutskever2013importance}, or
        AdamW~\citep{loshchilov2019decoupled}.
  \item \textbf{Schedule.} Linear warmup~\citep{goyal2017accurate} over 5\% of the steps (at
        most 500), then cosine decay~\citep{loshchilov2017sgdr} to 5\% of the peak rate.
        This differs from \katago{}'s continuous fixed rate, because v4 (like v3) trains a
        fresh fine-tuning run each cycle. \Cref{sec:limitations} discusses that choice.
  \item \textbf{Stability.} Gradients are clipped at norm 5.
  \item \textbf{Mixed precision.} On CUDA the trainer uses bf16 autocast where supported, or
        fp16 with loss scaling~\citep{micikevicius2018mixed}. \code{torch.compile} and
        channels-last layout are opt-in and limited to CUDA. On the Windows CPU we tested,
        Inductor compilation did not finish in 10 minutes.
\end{itemize}

\paragraph{Exponential moving average.}
The trainer keeps a shadow copy $\bar\theta$ of all parameters and batch-norm statistics,
updated after every step:
\begin{equation}
  \bar\theta \leftarrow d_t\,\bar\theta + (1-d_t)\,\theta,\qquad
  d_t = \min\!\left(d,\ \frac{1+t}{10+t}\right),\quad d = 0.999 .
\end{equation}
Evaluation, best-checkpoint selection and export all use $\bar\theta$. This is iterate
averaging in the sense of \citet{polyak1992acceleration}. It plays the same role as
\katago{}'s EMA of weight snapshots (decay 0.75 over snapshots taken every 250k samples)
and relates to SWA~\citep{izmailov2018averaging}. Our averaging is per step and much
shorter (horizon ${\approx}1/(1-d)=1000$ steps), which suits cycles of a few thousand
steps. The warm-up term keeps random early weights from lingering in a short run. The
update is one fused \code{\_foreach\_lerp\_} over all floating-point tensors.

\paragraph{Checkpoints and stopping.}
Validation runs every epoch or every $k$ steps, and early stopping has a patience counted
in evaluations. Every write is atomic (temporary file, then rename):
\begin{itemize}
  \item \code{best.pt} holds the best EMA weights as a plain state dict, so every v3
        consumer still works.
  \item \code{best.ckpt} adds the architecture and the metrics.
  \item \code{last.pt} holds the full state: weights, EMA, optimiser, scheduler, loss
        scaler and counters.
\end{itemize}
\code{--continue} restores that full state and resumes at the same step. A preempted
spot instance therefore loses at most \code{ckpt\_every} steps of training
(\cref{sec:infra}).

\subsection{Architectures}\label{sec:learner-arch}

Two families share one interface. Both return policy logits, value and ownership, and the
newer one also returns score and reply-policy outputs.

\paragraph{Legacy.}
This is the v3 network, bit-compatible with the champion's state dict:
\begin{itemize}
  \item A convolutional stem and $b$ post-activation residual blocks with $c$ channels.
  \item The flattened trunk is concatenated with an embedding of the global inputs.
  \item Fully connected policy and value heads read that vector, plus a $1\times1$
        ownership convolution.
\end{itemize}
The fully connected heads dominate the model. For 6$\times$64 they hold 766{,}484 of its
1{,}215{,}444 parameters (63\%). Their size also grows with $S^2$, which ties the network
to one board size.

\paragraph{\code{gpool}.}
This family follows \katago{}'s design:
\begin{itemize}
  \item \textbf{Trunk.} Pre-activation residual blocks~\citep{he2016identity}. The global
        inputs enter as a per-channel bias on the stem.
  \item \textbf{Global pooling.} Every third block routes $c_g$ of its channels after the
        first convolution through BN, ReLU and mean-and-max pooling, then a linear layer
        whose output biases the remaining $c-c_g$ channels. This is \katago{}'s
        global-pooling bias, without the board-width-scaled mean. That term is a constant
        multiple of the mean on a fixed $9\times9$ board and would need to be added back
        for mixed sizes.
  \item \textbf{Initialisation.} The second convolution of each block starts at zero, so
        every block begins as the identity. This is similar in spirit to
        \citet{goyal2017accurate}'s zero-$\gamma$ initialisation. It targets the seed
        sensitivity recorded in the repository's ADR~0005, where one 6$\times$32 seed
        diverged. We have not re-run that seed study.
  \item \textbf{Policy head.} A $1\times1$ convolution with a pooled bias produces two
        maps: the move policy and the opponent-reply policy. Two pass logits come from the
        pooled features.
  \item \textbf{Value head.} It pools a $1\times1$ convolution and feeds a two-layer MLP
        producing value ($\tanh$) and score. Ownership comes from the value head's spatial
        features, as in \katago{}.
\end{itemize}
No layer depends on $S$, so the parameter count is the same for every board size. A test
runs one network on $9\times9$, $13\times13$ and $19\times19$ inputs.

\Cref{tab:learner-arch} lists the presets. One caveat matters. \code{gpool}-6$\times$64 has
$2.7\times$ fewer parameters than legacy 6$\times$64, because the fully connected heads
are gone, but it needs about the \emph{same} multiply-accumulates per position. Its CPU
latency is no lower (\cref{sec:eval}). The parameter saving is not a speed saving.

\begin{table}[htbp]
\centering
\small
\begin{tabular}{@{}lrrrr@{}}
\toprule
Preset & Params & Head params & MACs / position & ONNX size \\
\midrule
legacy-6$\times$64 (v3 champion) & 1{,}215{,}444 & 766{,}484 & 37.0\,M & 4.9\,MB \\
legacy-4$\times$48 & 682{,}356 & 511{,}908 & 14.2\,M & 2.7\,MB \\
\code{gpool}-6$\times$64 & 446{,}631 & 12{,}935 & 35.2\,M & 1.3\,MB \\
\code{gpool}-10$\times$96 & 1{,}634{,}687 & 16{,}007 & 130.7\,M & 4.1\,MB \\
\code{gpool}-15$\times$128 & 4{,}321{,}575 & 33{,}223 & 345.7\,M & 10.2\,MB \\
\bottomrule
\end{tabular}
\caption{Architecture presets on $9\times9$. ``Head params'' counts every parameter
outside the stem and trunk. MACs count convolution and linear layers only, measured with
forward hooks on a single position. ONNX sizes are for policy+value exports.}
\label{tab:learner-arch}
\end{table}

\subsection{Changing architecture by distillation}\label{sec:learner-distill}

\katago{} grows its network during a run by training the next size concurrently on the
same data and switching once its loss catches up~\citep{wu2019accelerating}. That needs a
second trainer for the overlap period. v3 could not change size at all without
re-bootstrapping, estimated at 4--5 days (\cref{sec:bg-gofer}).

v4 adds knowledge distillation~\citep{hinton2015distilling} from the current champion. Any
checkpoint, of any architecture, can be passed as \code{--teacher}. On the same augmented
batch, the teacher's outputs add
\begin{equation}
  L_{\mathrm{distill}} = \mathrm{mean}_B\,\mathrm{CE}\!\left(\mathrm{softmax}(\ell^{T}_i),\,\hat\pi_i\right)
    + \mathrm{mean}_B\,(\hat v_i - v^{T}_i)^2
\end{equation}
to \cref{eq:loss}, with $w_d=1$ by default.

The student therefore learns from the replay targets and from the champion's own policy
and value on every position in the window, including fast-search positions whose policy
targets are masked. The added cost is one teacher forward pass per batch. The student must
still pass the ordinary gate against the champion before it is used for self-play. Whether
distillation actually avoids the re-bootstrap cost is an empirical question we have not
answered (\cref{sec:limitations}).

\subsection{Export}\label{sec:learner-export}

The exporter recovers the architecture in one of three ways:
\begin{itemize}
  \item from \code{best.ckpt} or \code{last.pt};
  \item from a \code{best.ckpt} stored next to a bare \code{best.pt};
  \item by inference from tensor shapes, for pre-v4 checkpoints.
\end{itemize}
Older checkpoints without the ownership head load with that head freshly initialised. Any
other missing or unexpected key is an error.

The graph is exported with ONNX opset 18~\citep{onnx2017} and exposes only
\code{policy\_logits} and \code{value}, the engine's contract. The architecture, a hash of
the weights, and the training step are embedded as ONNX metadata.

Every export then runs a parity check. ONNX Runtime and PyTorch evaluate 16 random
positions, and the export fails if any output differs by more than $10^{-3}$. The check
caught a real problem while being built. The legacy TorchScript exporter restores the
\emph{wrapper} module's training flag recursively after export, which silently switched
batch norm back to batch statistics in the reference. The reference output then differed
from the ONNX graph by 1.8. The fix was to put the wrapper, not just the inner network,
in evaluation mode.

\code{--quantize-int8} also writes a dynamically quantised copy~\citep{jacob2018quantization}
for CPU inference. It has not been gated for strength and is not used by default.
